\documentclass[a4paper,11pt]{article}
\usepackage{iftex}
\ifPDFTeX
\usepackage[T1]{fontenc}
\usepackage[utf8]{inputenc}
\usepackage{lmodern}
\else
\usepackage{fontspec}
\setmainfont{Latin Modern Roman}
\fi
\usepackage[ngerman]{babel}
\usepackage{amsmath,amssymb,amsthm,mathtools}
\usepackage{booktabs,tabularx,array}
\usepackage[a4paper,top=24mm,bottom=25mm,left=27mm,right=27mm]{geometry}
\usepackage{microtype}
\usepackage[hidelinks]{hyperref}
\newcommand{\Pp}[1]{\mathcal P_+(#1)}
\newcommand{\HL}{\mathcal H_L}
\newcommand{\HR}{\mathcal H_R}
\newcommand{\Sing}{\operatorname{Sing}}
\newtheorem{satz}{Satz}
\newtheorem{lemma}{Hilfssatz}
\theoremstyle{remark}
\newtheorem*{bemerkung}{Anmerkung}
\setlength{\parindent}{0pt}
\setlength{\parskip}{.55em plus .1em minus .1em}
\setlength{\emergencystretch}{1.2em}
\allowdisplaybreaks[2]
\hypersetup{pdftitle={Band 37 - Dreiprodukt-Rekonstruktion - Lesefassung},
 pdfauthor={Martin Kunze},pdfsubject={Endliche Halbgruppen mit drei Produktwerten und konstanten Nullprodukten bei drei Faktoren}}
\title{Wie Zeilen und Spalten die Halbgruppe bestimmen\\[.5em]
\large Rekonstruktion bei drei Produktwerten\\Lesefassung zu Band 37}
\author{Martin Kunze}
\date{5. Oktober 2026}
\begin{document}
\maketitle

Die Multiplikation nichtleerer Teilmengen enthält die Multiplikation der
Grundelemente als Einermengenkopie. Ein Isomorphismus dieser
Potenzhalbgruppen muss diese Kopie jedoch nicht erhalten. Deshalb suchen
wir eine andere Größe, aus der sich die Grundhalbgruppe zurückgewinnen
lässt: die Klassen von Elementen mit gleichen Zeilen und Spalten.

Diese Lesefassung entwickelt den vollständigen Beweis für eine bestimmte
endliche Klasse. Zwei Zählungen tragen den entscheidenden Schritt: erst
die Zählung innerhalb eines Zeilen- oder Spalten-Nullmusters, dann die
Zählung innerhalb des gemeinsamen Musters. Die zweite verhindert, dass
zwei verschiedene Begleiter bei der Rekonstruktion einen zusätzlichen
Punkt erzeugen.

Für eine Halbgruppe \(Q\) schreiben wir
\[
 Q^2=\{xy:x,y\in Q\},\qquad Q^3=\{xyz:x,y,z\in Q\}.
\]
Ihre Potenzhalbgruppe ist
\[
 \Pp Q=\{A\subseteq Q:A\ne\varnothing\},\qquad
 AB=\{ab:a\in A,\ b\in B\}.
\]
Insbesondere wird die leere Menge nicht als Element verwendet.

\begin{satz}\label{tp:main}
Sei \(S\) eine endliche Halbgruppe mit
\[
 S^2=\{0,u,v\},\qquad S^3=\{0\},
\]
wobei \(0,u,v\) paarweise verschieden sind. Dann gilt für jede
Halbgruppe \(T\):
\[
 \Pp S\cong\Pp T\quad\Longrightarrow\quad S\cong T.
\]
\end{satz}

Das Zeichen \(\cong\) verlangt die Existenz eines Isomorphismus.
Der Beweis konstruiert einen Grundisomorphismus aus einem gegebenen
Potenzisomorphismus; er behauptet keine Einermengentreue dieses
Potenzisomorphismus. Ausgangspunkt ist der am 5. Oktober 2026
bereitgestellte Beweisentwurf. Die vorliegende Fassung führt seine
Klasseneinteilungen und Zählargumente ausdrücklich aus.

\section{Die Zielhalbgruppe gehört zur gleichen Klasse}

Sei \(\Phi:\Pp S\to\Pp T\) ein Isomorphismus. Zunächst ist
\(T\) nichtleer und endlich: \(\Pp T\) ist eine nichtleere
endliche Menge, und \(t\mapsto\{t\}\) ist injektiv. Aus
\[
 2^{|S|}-1=|\Pp S|=|\Pp T|=2^{|T|}-1
\]
folgt \(|S|=|T|=:n\).

Jedes Dreifachprodukt in \(\Pp S\) ist \(\{0\}\). Daher ist
jedes Dreifachprodukt in \(\Pp T\) gleich
\(Z=\Phi(\{0\})\). Wählt man drei Einermengen auf der Zielseite,
so ist \(Z=\{abc\}\). Also \(Z=\{z\}\) und \(T^3=\{z\}\).
Jedes Vierfachprodukt ist auch ein Dreifachprodukt, etwa
\(abcd=(ab)cd\); deshalb \(T^4=\{z\}\). Schreibt man
\(z=abc\), so gilt für alle \(t\in T\)
\[
 zt=abct=z,\qquad tz=tabc=z.
\]
Damit ist \(z\) die Null von \(T\). Derselbe Schluss zeigt, dass
\(0\) die Null von \(S\) ist. Im Folgenden heißen beide Nullen
\(0\); insbesondere \(\Phi(\{0\})=\{0\}\).

Das Produktquadrat der Potenzhalbgruppe ist die Familie
\[
 (\Pp S)^2=\{AB:A,B\in\Pp S\}.
\]
Alle Mitglieder liegen in \(S^2\). Die folgenden sechs Mengen treten
stets als Produkte auf:
\[
 \{0\},\ \{u\},\ \{v\},\ \{0,u\},\ \{0,v\},\ \{0,u,v\}.
\]
Für \(\{u\}\) wählt man ein erzeugendes Paar \(ab=u\);
\(\{0,a\}\{b\}=\{0,u\}\) liefert seine Nullergänzung.
Entsprechend verfährt man für \(v\), und \(SS=S^2\) liefert die
Gesamtmenge. Nur \(\{u,v\}\) kann fehlen. Daher hat
\((\Pp S)^2\) sechs oder sieben Mitglieder. \(\Phi\) bildet dieses
Produktquadrat bijektiv auf \((\Pp T)^2\) ab.

Setze \(k=|T^2|\). Für \(k\le2\) gäbe es höchstens drei
Produktmengen. Für \(k\ge4\) gäbe es mindestens \(2k\): die
\(k\) Einermengen der Produktwerte, ihre \(k-1\) verschiedenen
Nullergänzungen und \(T^2=TT\). Diese letzte Menge ist von den
Ein- und Zweiermengen verschieden. Beide Alternativen widersprechen
den sechs oder sieben Produktmengen. Also \(|T^2|=3\).

Wir bezeichnen die beiden Nichtnullwerte auch auf der Zielseite mit
\(u,v\). Die Namen beziehen sich jeweils auf die betreffende
Halbgruppe. Da \(Q^3=\{0\}\) für \(Q=S,T\), gilt
\[
 Q^2Q=QQ^2=\{0\}.
\]
Alle Produktwerte haben somit vollständig nullige Zeilen und Spalten.

\section{Nullmuster und die Anzahl ihrer Grundelemente}

Für \(x\in Q\) definieren wir die Träger seiner Nichtnullprodukte:
\[
 L(x)=\{y\in Q:xy\ne0\},\qquad
 R(x)=\{y\in Q:yx\ne0\}.
\]
Für nichtleeres \(A\subseteq Q\) setzen wir
\(L(A)=\bigcup_{a\in A}L(a)\) und entsprechend
\(R(A)=\bigcup_{a\in A}R(a)\). Diese Mengen heißen linkes und
rechtes Nullmuster. Das gemeinsame Muster ist \((L(A),R(A))\).

Die Ordnungen dieser Muster sind intrinsisch erkennbar. Beispielsweise
\[
 L(A)\subseteq L(B)
 \ \Longleftrightarrow\ %
 \{X:BX=\{0\}\}\subseteq\{X:AX=\{0\}\}.
\]
Denn \(AX=\{0\}\) gilt genau dann, wenn \(X\) keinen Punkt
aus \(L(A)\) enthält. Die Rückrichtung der Äquivalenz prüft man
an Einermengen \(X\). Für \(R\) gilt der duale Zusammenhang.
Da \(\Phi\) die Null und alle Produktgleichheiten erhält, induziert
sie Ordnungsisomorphismen der realisierten linken, rechten und
gemeinsamen Muster.

\begin{lemma}\label{tp:null-count}
Unter diesen Musterzuordnungen bleiben die Anzahlen der Grundelemente
mit jedem einzelnen Muster erhalten.
\end{lemma}
\begin{proof}
Sei \(q\) ein realisiertes Muster und
\(D_q=\{x:\text{das Muster von }x\text{ ist höchstens }q\}\).
Weil Teilmengenmuster durch Vereinigung entstehen, hat eine nichtleere
Teilmenge genau dann Muster höchstens \(q\), wenn sie in \(D_q\)
liegt. Ihre Anzahl ist daher
\[
 d(q)=2^{N(q)}-1,\qquad N(q)=|D_q|=\log_2(d(q)+1).
\]
Die Ordnungszuordnung und \(\Phi\) erhalten \(d(q)\), also auch
\(N(q)\). Ist \(m(q)\) die Anzahl von Grundelementen mit genau
diesem Muster, so
\[
 m(q)=N(q)-\sum_{p<q}m(p).
\]
Induktion über die endliche Musterordnung erhält nun jede Zahl
\(m(q)\). Auch realisierte Teilmengenmuster ohne ursprüngliches
Grundelement sind zugelassen: Ihre Zahl \(m(q)\) ist Null.
\end{proof}

Diese Zählung funktioniert für einseitige Muster ebenso wie für das
gemeinsame Paar. Beide Formen werden im weiteren Beweis gebraucht.

\section{Der Fall erhaltener Produkt-Einermengen}

Wir sichern zuerst den Fall ab, in dem \(\Phi\) die drei
Produkt-Einermengen auf Produkt-Einermengen abbildet. Die folgende
Profilmethode erklärt, warum dafür keine Erhaltung aller übrigen
Einermengen erforderlich ist.

\begin{lemma}\label{tp:conditional}
Sind \(S,T\) endliche nichtleere Halbgruppen und ist
\(F:\Pp S\to\Pp T\) ein Isomorphismus mit
\[
 F[\Sing(S^2)]=\Sing(T^2),\qquad
 \Sing(I)=\{\{r\}:r\in I\},
\]
so gilt \(S\cong T\).
\end{lemma}
\begin{proof}
Definiere \(g:S^2\to T^2\) durch
\(F(\{r\})=\{g(r)\}\). Dies ist eine Bijektion. Wir vergleichen
nichtleere Teilmengen gemeinsam anhand aller konstanten Produkttests:
\[
 A\preceq B\ \Longleftrightarrow\
 \forall X\in\Pp S\ \forall r\in S^2:
 \begin{cases}
 BX=\{r\}\Rightarrow AX=\{r\},\\
 XB=\{r\}\Rightarrow XA=\{r\}.
 \end{cases}
\]
Das ist eine Präordnung. Gleiche Profile bedeuten
\(A\preceq B\) und \(B\preceq A\). Ihr Quotient ist eine
endliche teilweise geordnete Menge.

Für einen konstanten Produktwert gilt
\[
 AX=\{r\}\ \Longleftrightarrow\
 \forall a\in A:\ \{a\}X=\{r\}.
\]
Die Teilprodukte sind nichtleer; eine Vereinigung kann deshalb genau
dann eine Einermenge sein, wenn jedes Teilprodukt diese Einermenge
ist. Es folgt
\(A\preceq B\) genau dann, wenn alle \(\{a\}\preceq B\)
sind. Für ein Profil \(p\) ist die untere Profilfamilie somit die
volle nichtleere Potenzmenge von
\(D_p=\{s:[\{s\}]\le p\}\). Ihre Größe ist
\(2^{|D_p|}-1\).

\(F\) transportiert sämtliche Tests mit der Farbzuordnung \(g\).
Also bleiben die unteren Profilgrößen erhalten. Dieselbe endliche
Zählung wie in Hilfssatz~\ref{tp:null-count} liefert gleich große
Fasern von Grundelementen mit genau einem Profil.

Die Vorschriften \(r\mapsto g(r)\) für \(r\in S^2\) liegen
in den passenden Profilfasern, denn \(F(\{r\})=\{g(r)\}\).
Sie sind gemeinsam injektiv. In jeder Faser entfernt man die so
markierten Punkte und ihre Bilder; die gleich großen endlichen
Komplemente können bijektiv verbunden werden. Dies liefert eine
Bijektion \(f:S\to T\), die auf \(S^2\) mit \(g\)
übereinstimmt und für jedes \(s\) die Einermenge \(\{f(s)\}\)
profilgleich zu \(F(\{s\})\) macht.

Seien \(s,t\in S\), \(r=st\). Dann
\(F(\{s\})F(\{t\})=\{g(r)\}\). Gleiche Profile erlauben
bei jeder konstanten Produktgleichheit den Austausch erst des linken,
dann des rechten Faktors. Also
\[
 \{f(s)f(t)\}=\{g(r)\}=\{f(r)\}.
\]
Damit erhält \(f\) die Multiplikation und ist ein Isomorphismus.
\end{proof}

\section{Der einzige noch zu behandelnde Bildtyp}

Für \(A\in\Pp T\) und einen Produktwert \(Z\in(\Pp T)^2\)
schreiben wir
\[
 \ell_Z(A)=|\{B\in\Pp T:AB=Z\}|,\qquad
 \rho_Z(A)=|\{B\in\Pp T:BA=Z\}|.
\]
Für einen tatsächlichen Produkt-Singleton \(\{r\}\) ist
\(\ell_{\{r\}}(A)=2^k-1\): Die passenden rechten Faktoren
sind genau die nichtleeren Teilmengen der Spaltenmenge
\(\{t:A\{t\}=\{r\}\}\). Dasselbe gilt für \(\rho\).
Die Zahl Null ist dabei zugelassen, mit \(k=0\).

Eine Produktmenge \(Z\), die Null und einen Nichtnullwert enthält,
besitzt dagegen eine positive gerade Faktorisierungszahl für mindestens
einen linken Faktor. Wähle \(A_0B_0=Z\) und setze
\(A=A_0\cup\{0\}\). Dann \(AB_0=Z\).
Schreibe \(W=T^2=\{0,u,v\}\). Jeder rechte Faktor zerfällt
eindeutig in seinen Anteil außerhalb von \(W\) und eine Teilmenge
von \(W\). Ein Faktor, der mit \(A\) den Wert \(Z\) ergibt,
muss einen nichtleeren Anteil außerhalb von \(W\) besitzen.
Zu jedem passenden solchen Anteil sind sämtliche acht
\(W\)-Dekorationen erlaubt: \(W\) erzeugt nur Null, und die
Null ist wegen \(0\in A\) ohnehin bereits im Produkt vorhanden.
Also ist \(\ell_Z(A)\) positiv und durch acht teilbar.

Isomorphismen erhalten diese Faktorisierungszahlen für jeden festen
Faktor. Das Bild eines Nichtnull-Singletons kann daher keine solche
Nullergänzung sein. Die Bilder von \(\{u\}\) und \(\{v\}\)
liegen unter
\[
 \{u\},\quad\{v\},\quad\{u,v\}.
\]
Sind beide Bilder Einermengen, greift
Hilfssatz~\ref{tp:conditional}. Andernfalls wird genau eines der
Bilder \(\{u,v\}\). Durch getrennte Benennung der beiden
Nichtnullwerte auf Quell- und Zielseite dürfen wir also voraussetzen:
\begin{equation}\label{tp:normalform}
 \Phi(\{0\})=\{0\},\qquad
 \Phi(\{u\})=\{u\},\qquad
 \Phi(\{v\})=\{u,v\}.
\end{equation}
Hier liegt die eigentliche Schwierigkeit. Wir lösen sie im Folgenden
durch Zeilen- und Spaltenklassen, ohne diese Bildmengen in Einermengen
zurückverwandeln zu müssen.

\section{Gleiche Zeilen aus Faktorisierungszahlen}

Wir arbeiten jetzt im verbleibenden Fall \eqref{tp:normalform}:
\(\Phi(\{0\})=\{0\}\), \(\Phi(\{u\})=\{u\}\) und
\(\Phi(\{v\})=\{u,v\}\). Die beiden Grundhalbgruppen haben dieselbe
endliche Größe \(n\). Unter einer Zeile eines Elements \(a\) verstehen
wir die gesamte Funktion \(t\mapsto at\), unter seiner Spalte die
Funktion \(t\mapsto ta\). Zwei Zeilen sind also nur dann gleich, wenn
alle ihre Einträge übereinstimmen.

Für \(A\in\Pp{T}\) zählen wir die rechten Faktoren eines vorgegebenen
Produkts:
\[
 \ell_Z(A)=\bigl|\{B\in\Pp{T}:AB=Z\}\bigr|.
\]
Die drei Anzahlen zu \(Z=\{0\},\{u\},\{u,v\}\) sind sämtlich von
der Form \(2^k-1\), einschließlich \(k=0\). Um dies zu sehen, setzen
wir \(X=\Phi^{-1}(A)\). Die geeigneten Faktoren im Urbild sind genau
die nichtleeren Teilmengen der jeweiligen Mengen
\[
 E_r(X)=\{s\in S:X\{s\}=\{r\}\},\qquad r\in\{0,u,v\}.
\]
Hier wird wesentlich verwendet, dass die drei Urbildprodukte
Einermengen sind. Die Normalform gibt deshalb
\[
 \ell_{\{0\}}(A)=2^{|E_0(X)|}-1,\quad
 \ell_{\{u\}}(A)=2^{|E_u(X)|}-1,\quad
 \ell_{\{u,v\}}(A)=2^{|E_v(X)|}-1.
\]
Diese drei Mengen sind disjunkt. Damit ist die ganze Zahl
\begin{equation}
 d_L(A)=\log_2\bigl(\ell_{\{0\}}(A)+1\bigr)
       +\log_2\bigl(\ell_{\{u\}}(A)+1\bigr)
       +\log_2\bigl(\ell_{\{u,v\}}(A)+1\bigr)
       \le n
 \label{tp:middle-left-degree}
\end{equation}
definiert. Die Logarithmen dienen nur als Kurzschreibweise für die
eindeutig bestimmten endlichen Exponenten.

Gleichheit gilt genau dann, wenn für jedes \(s\in S\) das Produkt
\(X\{s\}\) eine Einermenge ist. Das bedeutet: In jeder Spalte haben
alle Elemente von \(X\) denselben Eintrag. Somit
\begin{equation}
 d_L(A)=n
 \quad\Longleftrightarrow\quad
 \text{alle Elemente von }\Phi^{-1}(A)\text{ haben dieselbe Zeile}.
 \label{tp:middle-homogeneous-source}
\end{equation}
Wir nennen die so bestimmten Mengen \(\HL\). Dual setzen wir
\[
 \rho_Z(A)=\bigl|\{B\in\Pp{T}:BA=Z\}\bigr|
\]
und definieren \(d_R(A)\) durch dieselbe Summe mit \(\rho\) anstelle
von \(\ell\). Die Familie \(\HR=\{A:d_R(A)=n\}\) besteht genau
aus den Bildern von Mengen gleicher ursprünglicher Spalten.

Um diese Familien in Klassen zu gliedern, verwenden wir die
\emph{vollständigen Translationen}:
\[
 A\equiv_L B\ \Longleftrightarrow\
       \forall Y\in\Pp{T}\ (AY=BY),\qquad
 A\equiv_R B\ \Longleftrightarrow\
       \forall Y\in\Pp{T}\ (YA=YB).
\]
Die erste Äquivalenz wird auf \(\HL\), die zweite auf \(\HR\)
betrachtet. \(\Phi\) erhält und reflektiert beide, weil sie alle
Testmengen bijektiv zuordnet und die Produkte erhält.

Im Urbild besteht jede linke Klasse aus genau den nichtleeren
Teilmengen einer Klasse gleicher Zeilen in \(S\). Verschiedene
Zeilen werden bereits durch eine passende Singleton-Testmenge
unterschieden. Eine Klasse von \(k\) ursprünglichen Zeilen liefert
daher eine Familie mit \(2^k-1\) Mitgliedern. Dasselbe gilt rechts
für gleiche Spalten. Hier werden vollständige Zeilen verglichen;
gleiche Nullmuster allein wären nicht ausreichend.

\subsection{Die möglichen Mersenne-Anzahlen}

Zunächst lösen wir die kleine Zahlenfrage, die in den Mengenprodukten
wiederkehrt. Für nichtnegative ganze Zahlen \(p,q,r\) sei
\[
 M=2^{p+q+r}-2^p-2^q+1.
\]
Ist \(M>0\) von der Form \(2^k-1\), so gilt genau eine der
beiden Möglichkeiten
\begin{equation}
 p=q=0,\ r>0;\qquad
 r=0,\ p,q>0,\ \min(p,q)=1.
 \label{tp:middle-mersenne-cases}
\end{equation}
Im ersten Fall ist \(k=r\), im zweiten \(k=\max(p,q)\).

Für \(p=q=0\) ist die Behauptung unmittelbar. Ist genau einer von
\(p,q\) positiv, etwa \(p>0,q=0\), dann ist
\(M=2^p(2^r-1)\). Diese Zahl ist entweder Null oder positiv und
gerade, also keine positive Zahl \(2^k-1\). Bleibt \(p,q>0\).
Wären beide mindestens zwei, so wäre
\[
 M+1=2^{p+q+r}-2^p-2^q+2\equiv2\pmod4.
\]
Eine Zweierpotenz mit diesem Rest ist nur \(2\), während hier
\(M\ge9\) gilt. Also ist einer der beiden Exponenten eins.
Für \(q=1\) erhalten wir
\[
 M+1=2^p\bigl(2^{r+1}-1\bigr).
\]
Der zweite Faktor ist ungerade und positiv. Das Produkt ist genau
dann eine Zweierpotenz, wenn dieser Faktor eins ist, also \(r=0\).
Damit ist \(k=p\). Der Fall \(p=1\) ist symmetrisch. Die
angegebenen Möglichkeiten liefern umgekehrt tatsächlich die
behaupteten Anzahlen.

Nun betrachten wir eine einzelne Zeile von \(T\), mit \(z\) Nullen,
\(p>0\) Einträgen \(u\) und \(q>0\) Einträgen \(v\). Ihre
Faktorisierungszahlen zu den drei ausgezeichneten Produkten sind
\[
 2^z-1,\qquad 2^p-1,\qquad(2^p-1)(2^q-1).
\]
Die letzte Zahl ist die vorige Formel mit \(r=0\). Also muss
\(\min(p,q)=1\) sein. Wäre \(q=1<p\), so ergäbe
\eqref{tp:middle-left-degree}
\[
 d_L(\{a\})=z+2p>z+p+1=n,
\]
ein Widerspruch. Folglich enthält jede Zeile mit beiden Nichtnullwerten
\emph{genau einen Eintrag \(u\)}. Mit \(d_R\) folgt dieselbe
Aussage für jede Spalte mit beiden Nichtnullwerten.

\subsection{Welche Zielmengen zu gleichen ursprünglichen Zeilen gehören}

Für eine beliebige nichtleere Menge \(A\subseteq T\) teilen wir die
Spalten nach ihren Produkten \(A\{t\}\) ein. Es seien \(z,p,q,r\)
die Anzahlen der Typen \(\{0\},\{u\},\{v\},\{u,v\}\).
Weitere Spalten können die Typen \(\{0,u\},\{0,v\},\{0,u,v\}\)
besitzen; ihre gemeinsame Anzahl heiße \(h\). Damit
\(n=z+p+q+r+h\).

Ein Produkt gleich \(\{u,v\}\) kann keine Spalte mit Nullanteil
verwenden. Unter den übrigen Spalten darf man alle Auswahlen nehmen,
die gemeinsam beide Farben liefern. Man zieht die ausschließlich
\(u\)-farbenen und ausschließlich \(v\)-farbenen Auswahlen ab.
Dies ergibt
\begin{equation}
 \ell_{\{u,v\}}(A)=2^{p+q+r}-2^p-2^q+1,\qquad
 \ell_{\{0\}}(A)=2^z-1,\qquad
 \ell_{\{u\}}(A)=2^p-1.
 \label{tp:middle-rectangle-count}
\end{equation}

Ist die erste Zahl Null und \(A\in\HL\), so ist
\(d_L(A)=z+p=n\). Folglich sind \(q=r=h=0\); alle Elemente
von \(A\) besitzen dieselbe Zeile mit Einträgen nur \(0,u\).
Dies schließt die vollständig nullige Zeile ausdrücklich ein.

Ist die erste Zahl positiv, verwenden wir
\eqref{tp:middle-mersenne-cases}. Für \(p=q=0,r>0\) ist
\(d_L(A)=z+r\). Die Gleichheit mit \(n\) erzwingt \(h=0\).
Es gibt daher eine nichtleere Menge \(D\subseteq T\) mit
\begin{equation}
 A\{t\}=\{u,v\}\quad(t\in D),\qquad
 A\{t\}=\{0\}\quad(t\notin D).
 \label{tp:middle-mixed-rows}
\end{equation}
Im anderen Fall sind \(r=0,p,q>0\). Jedes Element von \(A\)
besitzt dann mindestens \(p\) Einträge \(u\) und mindestens
\(q\) Einträge \(v\). Die Zeileneinschränkung erzwingt \(p=1\).
Nun ist \(d_L(A)=z+1+q\), sodass wieder \(h=0\) folgt. Alle
Spaltenprodukte sind Einermengen, also sind alle Zeilen in \(A\)
gleich.

Damit ist \(\HL\) vollständig beschrieben: Es enthält die
nichtleeren Mengen gleicher Zeilen, sofern diese gemeinsame Zeile
kein nichtnulliger reiner \(v\)-Typ ist, und zusätzlich die Mengen
\eqref{tp:middle-mixed-rows}. Umgekehrt erfüllen alle diese Mengen
\(d_L=n\). Eine reine \(v\)-Zeile mit nichtleerem Träger hat
dagegen \(d_L=z<n\).

\subsection{Wie viele Mengen die zusätzliche Zeilenklasse besitzen muss}

Fixiere ein nichtleeres Zeilen-Nullmuster \(D\) und setze
\[
 C_D=\{a\in T:\{t:at\ne0\}=D\},\qquad
 V_D=\{a\in C_D:\forall t\in D\ (at=v)\},\qquad b=|V_D|.
\]
Die gewöhnlichen Klassen gleicher Zeilen in \(C_D\), außer
\(V_D\), erscheinen als eigene \(\equiv_L\)-Klassen in
\(\HL\). Eine solche Punktklasse \(R\) liefert genau
\(\Pp{R}\), also \(2^{|R|}-1\) Mengen. Ihre im Urbild
rekonstruierte Punktklasse besitzt deshalb ebenfalls \(|R|\)
Elemente.

Daneben gibt es höchstens eine zusätzliche Klasse \(\mathcal G_D\):
die Mengen aus \eqref{tp:middle-mixed-rows}. Diese besitzen dieselbe
vollständige linke Translation, denn die Produkte mit jeder
Singleton-Testmenge sind festgelegt. Sie können keiner gewöhnlichen
Klasse angehören, da ihre Produkte auf \(D\) Zweiermengen sind.
Unterschiedliche \(D\) liefern ebenfalls unterschiedliche
Translationen. Jedes ihrer Mitglieder liegt in \(\Pp{C_D}\).

Sei \(g=|\mathcal G_D|\). Ist die Familie nichtleer, so hat sie
im Urbild die Größe \(2^k-1\) für eine ursprüngliche Zeilenklasse;
bei leerer Familie setzen wir \(k=0\). Die Nullmusterzählung aus
Hilfssatz~\ref{tp:null-count} erhält die Anzahl ursprünglicher Elemente
in diesem Muster. Die gewöhnlichen Klassen rekonstruieren zusammen
\(|C_D|-b\) Elemente, die zusätzliche Klasse rekonstruiert \(k\).
Daher
\begin{equation}
 (|C_D|-b)+k=|C_D|,\qquad
 k=b,\qquad g=2^b-1.
 \label{tp:middle-extra-count}
\end{equation}
Insbesondere ist die zusätzliche Klasse genau dann vorhanden,
wenn \(V_D\) nichtleer ist.

\subsection{Die Begleiter reiner \texorpdfstring{\(v\)}{v}-Zeilen}

Sei nun \(b>0\). Wegen \eqref{tp:middle-extra-count} existiert eine
Menge aus \(\mathcal G_D\). Für jedes \(t\in D\) liefert sie
mindestens einen Eintrag \(u\), während \(V_D\) einen Eintrag
\(v\) liefert. Diese Spalte enthält beide Farben und deshalb
\emph{global genau ein \(u\)}. Sein Zeilenelement heiße \(p_t\).
Es liegt in \(C_D\). Setze \(U=\{p_t:t\in D\}\).

Ein Element von \(C_D\setminus V_D\) besitzt irgendwo einen
Eintrag \(u\), muss also eines der \(p_t\) sein. Folglich
\(C_D=U\mathbin{\dot\cup}V_D\). Wäre \(|U|\ge2\), könnte
keine Zeile aus \(U\) auf ganz \(D\) ausschließlich \(u\)
enthalten: Ein solches Element wäre in jeder Spalte das einzige
\(u\), und alle \(p_t\) wären dasselbe Element. Somit besäße
jede Zeile aus \(U\) beide Farben und genau ein \(u\).

In jeder Spalte aus \(D\) liefert nun \(p_t\) das \(u\), jedes
andere Element von \(U\) das \(v\). Also wäre bereits \(U\)
eine Menge aus \(\mathcal G_D\). Jedes Mitglied dieser Klasse muss
alle \(p_t\) enthalten, da diese allein das benötigte \(u\)
liefern. Die Mitglieder wären daher genau \(U\cup B\),
\(B\subseteq V_D\), einschließlich \(B=\varnothing\).
Ihre Anzahl wäre \(2^b\), entgegen
\eqref{tp:middle-extra-count}.

Also ist \(U=\{p_D\}\). Dieses Element hat dasselbe Nullmuster
wie \(V_D\) und ausschließlich \(u\) als Nichtnullwerte. Seine
ganze Klasse gleicher Zeilen ist ein Singleton: Jede gleiche Zeile
läge in \(C_D\setminus V_D=\{p_D\}\). Genau die Mengen
\begin{equation}
 \mathcal G_D=\bigl\{\{p_D\}\cup B:
                 \varnothing\ne B\subseteq V_D\bigr\}
 \label{tp:middle-row-companion}
\end{equation}
besitzen die zusätzliche Translation. Der Begleiter muss das
\(u\) liefern, mindestens ein Element aus \(V_D\) das \(v\).

Alle Schritte gelten dual für Spalten. Zu jeder nichtnulligen Klasse
reiner \(v\)-Spalten mit Nullmuster \(D\) gehört ein eindeutiger
Begleiter \(q_D\), dessen einzige Nichtnullwerte \(u\) sind und
dessen ganze Spaltenklasse \(\{q_D\}\) ist. Die zusätzliche
\(\equiv_R\)-Klasse in \(\HR\) besteht genau aus
\(\{q_D\}\cup B\) für nichtleere Teilmengen der betreffenden
reinen \(v\)-Spaltenklasse. Damit kennen wir beide Arten
rekonstruierter Klassen; als Nächstes müssen ihre Schnittpunkte
zusammenpassen.


Wir halten das soeben bewiesene Ergebnis fest.
\begin{lemma}\label{tp:companions}
Im verbleibenden Normalfall \eqref{tp:normalform} besitzt jede nichtnullige
Klasse $R$ reiner $v$-Zeilen einen eindeutigen
Begleiter $p$ mit demselben Nullmuster und ausschließlich $u$ als
Nichtnullwerten. Seine ganze Zeilenklasse ist $\{p\}$, und die zugehörige
zusätzliche Klasse in $\HL$ besteht genau aus $\{p\}\cup B$ mit
$\varnothing\ne B\subseteq R$. Dual gilt die entsprechende Aussage
für reine $v$-Spalten und ihre Begleiter.
\end{lemma}

\section{Zeilen und Spalten zusammenführen}
\label{tp:joint-classes}

Die bisherige Zählung behandelt gleiche Zeilen und gleiche Spalten
getrennt. Für eine Bijektion der ursprünglichen Elemente müssen wir nun
ihre Schnittmengen betrachten. Genau hier wird ausgeschlossen, dass zwei
verschiedene Begleiter gemeinsam ein zusätzliches Element vortäuschen.

Sei $R$ eine Klasse gleicher Zeilen in $T$, $K$ eine Klasse gleicher
Spalten und $E=R\cap K$. Bei $E\ne\varnothing$ ist $E$ eine Klasse
gleicher Zeilen \emph{und} Spalten. Ein \emph{$v$-Typ} bezeichnet im
Folgenden eine nichtnullige Zeile beziehungsweise Spalte, deren sämtliche
Nichtnulleinträge $v$ sind. Eine vollständig nullige Zeile oder Spalte ist
kein $v$-Typ.

Zu jeder Zeilenklasse $R$ gehört nach dem Begleiterlemma
\ref{tp:companions} genau eine Klasse vollständiger linker Translationen
innerhalb von $\HL$. Wir schreiben dafür $\mathcal F_R$. Ist $R$ kein
$v$-Typ, so ist $\mathcal F_R=\Pp{R}$. Ist $R$ ein $v$-Typ mit
Begleiter $p$, so gilt
\[
 \mathcal F_R=
 \bigl\{\{p\}\cup B:\varnothing\ne B\subseteq R\bigr\}.
\]
Insbesondere ist die ganze ursprüngliche Zeilenklasse des Begleiters
$p$ die Einermenge $\{p\}$. Entsprechend gehört zu jeder Spaltenklasse
$K$ eine Klasse $\mathcal G_K$ vollständiger rechter Translationen in
$\HR$. Sie ist $\Pp{K}$, wenn $K$ kein $v$-Typ ist, und andernfalls
\[
 \mathcal G_K=
 \bigl\{\{q\}\cup B:\varnothing\ne B\subseteq K\bigr\},
\]
wobei $q$ der Spaltenbegleiter ist und seine ganze ursprüngliche
Spaltenklasse $\{q\}$ beträgt. Diese Familien werden nach vollständigen
Translationen unterschieden; das Nullmuster allein definiert ihre Klassen
nicht.

\subsection{Die fünf möglichen Schnittmengen}

Wir bestimmen $\mathcal F_R\cap\mathcal G_K$ unmittelbar aus diesen
Beschreibungen. Wenn keine der beiden Klassen ein $v$-Typ ist, erhalten
wir einfach die nichtleeren Teilmengen von $E$. Ist nur $R$ ein $v$-Typ,
so muss sein Begleiter $p$ zusätzlich in $K$ liegen. Die Mengen sind dann
$\{p\}\cup B$ mit $\varnothing\ne B\subseteq E$. Der Fall, dass nur
$K$ ein $v$-Typ ist, verläuft entsprechend mit $q\in R$.

Sind beide Klassen $v$-Typen und stimmen ihre Begleiter überein, so
werden wieder genau die nichtleeren Teilmengen von $E$ um den gemeinsamen
Begleiter ergänzt. Sind dagegen $p\ne q$, muss $q$ zum Zeilenanteil und
$p$ zum Spaltenanteil gehören. Deshalb sind $q\in R$ und $p\in K$
notwendig. Unter diesen Voraussetzungen lautet die Familie genau
\begin{equation}
 \mathcal F_R\cap\mathcal G_K
 =\bigl\{\{p,q\}\cup B:B\subseteq E\bigr\}.
 \label{tp:crossing-family}
\end{equation}
Hier ist auch $B=\varnothing$ erlaubt: Die beiden erforderlichen
nichtleeren Anteile enthalten bereits $q$ beziehungsweise $p$.

Damit ergibt sich die vollständige Tabelle. Fehlt eine dort genannte
Voraussetzung, ist die betreffende Familie leer.
\begin{center}
\small
\begin{tabularx}{\linewidth}{@{}
  >{\raggedright\arraybackslash}X
  >{\raggedright\arraybackslash}X
  >{\raggedleft\arraybackslash}p{0.20\linewidth}@{}}
\toprule
Typen von $R,K$ & Voraussetzung & Anzahl \\
\midrule
Keiner ein $v$-Typ & keine & $2^{|E|}-1$ \\
Nur $R$ ein $v$-Typ & $p\in K$ & $2^{|E|}-1$ \\
Nur $K$ ein $v$-Typ & $q\in R$ & $2^{|E|}-1$ \\
Beide, $p=q$ & keine & $2^{|E|}-1$ \\
Beide, $p\ne q$ & $p\in K$, $q\in R$ & $2^{|E|}$ \\
\bottomrule
\end{tabularx}
\end{center}

Auf der Quellseite sind die entsprechenden Schnittmengen besonders
einfach. Eine Klasse gleicher Zeilen und Spalten in $S$ mit $k$ Elementen
liefert genau ihre $2^k-1$ nichtleeren Teilmengen. Der Isomorphismus
$\Phi$ transportiert sowohl die vollständigen linken als auch die
vollständigen rechten Translationen. Jede nichtleere Familie der Tabelle
ist daher das Bild einer solchen Quellfamilie. Ihre Größe muss $2^k-1$
mit $k\ge1$ sein.

Für den letzten Tabellenfall folgt somit
\[
 2^{|E|}=2^k-1.
\]
Die rechte Seite ist ungerade. Also muss $|E|=0$ und folglich $k=1$
gelten. Ein solcher Schnitt hätte auf der Zielseite keine ursprünglichen
Elemente in $E$, würde aber auf der Quellseite ein Element rekonstruieren.
Wir zeigen jetzt, dass auch diese verbleibende Möglichkeit ausgeschlossen
ist.

\subsection{Warum verschiedene Begleiter nicht kreuzen können}

Nehmen wir einen letzten Tabellenfall mit $p\ne q$ an. Die beiden
Begleiter haben dasselbe gemeinsame Nullmuster. Tatsächlich liegt $q$ in
$R$, während $p$ als Zeilenbegleiter dasselbe linke Nullmuster wie $R$
besitzt. Daher ist $L(p)=L(q)$. Ebenso liegt $p$ in $K$, während $q$
als Spaltenbegleiter dessen rechtes Nullmuster besitzt. Also gilt auch
$R(p)=R(q)$.

Innerhalb dieses gemeinsamen Nullmusters gehen zugleich zwei tatsächliche
Elemente bei der Rekonstruktion verloren. Die Zeilenklasse von $p$ ist
$\{p\}$ und kein $v$-Typ; seine Spaltenklasse ist $K$. Deshalb lautet
die zugehörige linke Familie $\mathcal F_{\{p\}}=\{\{p\}\}$.
Jede Menge aus $\mathcal G_K$ enthält jedoch den anderen Begleiter $q$.
Der Schnitt ist leer. Die wirkliche gemeinsame Klasse $\{p\}$ trägt
also kein rekonstruiertes Element bei. Dual ist die Spaltenklasse von
$q$ die Einermenge $\{q\}$, während seine Zeilenklasse $R$ ist. Jede
Menge aus $\mathcal F_R$ enthält $p$, sodass auch die gemeinsame Klasse
$\{q\}$ keinen Beitrag liefert.

Es kann innerhalb desselben gemeinsamen Nullmusters höchstens eine
zusätzliche Klasse dieser Art geben. Das linke Nullmuster legt eine
$v$-Zeile vollständig fest: Auf ihrem Nichtnullbereich stehen überall
$v$, außerhalb überall $0$. Deshalb gibt es dafür nur eine mögliche
Zeilenklasse $R$ vom $v$-Typ. Dasselbe gilt für die mögliche
Spaltenklasse $K$. Das zusätzliche Klassenpaar $(R,K)$ ist somit
eindeutig.

Die folgende Zählung bleibt dabei innerhalb eines festen gemeinsamen
Nullmusters. Jede nichtleere Schnittfamilie hat das linke Muster ihrer
Zeilenklasse $R$ und das rechte Muster ihrer Spaltenklasse $K$.
Hinzugefügte Begleiter ändern diese Muster nicht: Ein Zeilenbegleiter
hat dasselbe linke Muster wie $R$ und gehört nach der Tabellenbedingung
zur Spaltenklasse $K$; für einen Spaltenbegleiter gilt das dual.

Alle übrigen Schnittmengen der Tabelle liefern höchstens so viele
rekonstruierte Elemente, wie ihre tatsächliche Klasse $E$ enthält.
Eine Familie mit Größe $2^{|E|}-1$ liefert genau $|E|$ Elemente; eine
leere Familie liefert keines. Der angenommene Kreuzungsfall könnte
höchstens ein Element hinzufügen, nachdem mindestens die beiden
Elemente $p$ und $q$ verloren gingen. Die rekonstruierte Gesamtzahl im
betrachteten gemeinsamen Nullmuster wäre daher mindestens um eins
kleiner als die ursprüngliche Gesamtzahl.

Das widerspricht der Nullmusterzählung aus Hilfssatz~\ref{tp:null-count}, die diese
Gesamtzahlen unter $\Phi$ erhält. Der Fall verschiedener kreuzender
Begleiter tritt folglich nicht auf.

Nun bleiben ausschließlich die ersten vier Tabellenfälle. Auch eine
leere tatsächliche Klasse $E$ kann dort keine zusätzliche Familie
erzeugen. Für jedes gemeinsame Nullmuster ist daher die rekonstruierte
Punktzahl höchstens die tatsächliche Punktzahl. Da beide nach
Hilfssatz~\ref{tp:null-count} gleich sind, darf keine nichtleere tatsächliche
Klasse verloren gehen: Jede ihrer Schnittfamilien ist nichtleer und
liefert genau $|E|$ Elemente.

Damit sind die gemeinsamen Zeilen- und Spaltenklassen von $S$ bijektiv
den entsprechenden tatsächlichen Klassen von $T$ zugeordnet. Zugeordnete
Klassen haben dieselbe Größe. Außerdem enthält jede Menge ihrer
Schnittfamilie eine \emph{nichtleere} Teilmenge von $E$; gegebenenfalls
kommt der eine gemeinsame Begleiter hinzu. Diese zusätzliche Aussage
wird bei der Multiplikation entscheidend.

\section{Die rekonstruierte Bijektion erhält die Produkte}

Wähle in jeder zugeordneten Klassenpaarung eine Bijektion. Zusammen
ergeben diese eine Bijektion $f:S\to T$. Die drei Produktwerte $0,u,v$
liegen in beiden Halbgruppen in der Klasse vollständig nulliger Zeilen
und Spalten, weil $S^3=T^3=\{0\}$ gilt. Diese Klassen werden einander
zugeordnet. Deshalb können wir ihre Klassenbijektion so wählen, dass
\[
 f(0)=0,\qquad f(u)=u,\qquad f(v)=v.
\]
Weitere beidseitige Annihilatoren in dieser Klasse werden durch eine
beliebige Bijektion der verbleibenden Elemente zugeordnet.

Seien $s,t\in S$ und setze $A=\Phi(\{s\})$, $B=\Phi(\{t\})$.
Die zugeordnete tatsächliche Klasse von $f(s)$ trifft $A$; entsprechend
trifft die Klasse von $f(t)$ die Menge $B$. Wähle daraus $a\in A$ und
$b\in B$. Dann haben $a,f(s)$ dieselbe Zeile, während $b,f(t)$ dieselbe
Spalte haben. Folglich
\begin{equation}
 f(s)f(t)=a f(t)=ab\in AB.
 \label{tp:membership}
\end{equation}
Hier müssen $f(s)$ und $f(t)$ selbst nicht in ihren Potenzbildern liegen.
Die Gleichheit ihrer Zeilen beziehungsweise Spalten mit geeigneten
Elementen der Bilder genügt.

Ist $st=0$ oder $st=u$, so ist $AB$ nach
\eqref{tp:normalform} die entsprechende Einermenge. Aus
\eqref{tp:membership} folgt unmittelbar
$f(s)f(t)=st=f(st)$.

Sei schließlich $st=v$. Jetzt ist $AB=\{u,v\}$, sodass
\eqref{tp:membership} zunächst nur $f(s)f(t)\in\{u,v\}$ liefert.
Angenommen, dieses Produkt wäre $u$. Dann kann die Zeile von $f(s)$
kein $v$-Typ sein; sie besitzt mindestens diesen einen $u$-Eintrag.
Entsprechend kann die Spalte von $f(t)$ kein $v$-Typ sein. Für eine
gewöhnliche Zeilenklasse besteht die zugehörige Familie ausschließlich
aus ihren nichtleeren Teilmengen. Daher haben alle Elemente von $A$
dieselbe Zeile wie $f(s)$. Ebenso haben alle Elemente von $B$ dieselbe
Spalte wie $f(t)$.

Für jedes $a'\in A$, $b'\in B$ folgt deshalb
\[
 a'b'=f(s)b'=f(s)f(t)=u.
\]
Das würde $AB=\{u\}$ ergeben, im Widerspruch zu
$AB=\{u,v\}$. Also ist $f(s)f(t)=v=f(st)$. Damit erhält $f$ jedes
Produkt und ist der gesuchte Halbgruppenisomorphismus.

Der zuvor behandelte Fall, dass die beiden Nichtnullprodukt-Singletons
auf Singletons abgebildet werden, sowie der jetzt bewiesene Normalfall
\eqref{tp:normalform} erschöpfen alle Möglichkeiten. Somit ist die
Rekonstruktionsaussage für endliche Halbgruppen mit
$S^2=\{0,u,v\}$ und $S^3=\{0\}$ bewiesen. Die allgemeine endliche
Potenzhalbgruppenvermutung umfasst weitere Halbgruppen und wird dadurch
nicht entschieden.


\section*{Einordnung und Herkunft}

Der Satz behandelt genau endliche Halbgruppen mit drei Produktwerten und
konstanten Nullprodukten bei drei Faktoren. Das Ziel darf zunächst eine
beliebige Halbgruppe sein; seine Endlichkeit und seine Klassenmerkmale
werden im Beweis zurückgewonnen. Es wird weder verlangt, dass der gegebene
Potenzisomorphismus Einermengen erhält, noch dass er durch den konstruierten
Grundisomorphismus induziert wird.

Der bereitgestellte Entwurf wurde für diese Lesefassung mathematisch
gegengeprüft. Die Darstellung ergänzt insbesondere die positive gerade
Faktorisierungszahl, die Klassen vollständiger Translationen, die
Mersenne-Fallanalyse und die Zählung der gemeinsamen Nullmuster.
Eine Aussage über die Priorität des Satzes in der Literatur wird damit
nicht verbunden.

Die allgemeinen Bausteine stehen im Skript in Band~20
(endliche Faserzählung), Band~28 (Potenzhalbgruppen), Band~33
(Nullproduktprofile) und Band~37 (endliche Halbgruppen). Die frühere
Lesefassung \emph{Was Nullprodukte verraten} erläutert den
Zweiproduktfall. Der vorliegende Beweis erweitert die Zählidee auf den
hier behandelten Dreiproduktfall und benötigt dafür zusätzlich die
Vereinbarkeit von Zeilen- und Spaltenbegleitern.

Als endliche Kontrolle wurden die Zwischenbehauptungen an allen
15 Fünfelementklassen und allen 1\,622 Sechselementklassen dieser Klasse
geprüft. Die Gültigkeit für beliebige endliche Größe folgt aus dem oben
gegebenen Zählbeweis. Die allgemeine endliche Potenzhalbgruppenvermutung
umfasst weitere Klassen und wird hier nicht als vollständig gelöst
behauptet.

\end{document}
